\chapter{Software Verification, Numerical Verification, and Validation}
\label{ch:verification-and-validation}

\textbf{Status: evidence framework and retained-status interpretation; no new
scientific validation or uncertainty-quantification result is claimed.}

Scientific software can run successfully, reproduce its own output, and still
answer the wrong mathematical or physical question. This chapter separates the
evidence classes needed by the reduction program and states what each class can
and cannot establish.

\section{Claims precede evidence}

An evidence record begins with a bounded claim. The claim identifies the system,
quantity, representation, parameter domain, metric, tolerance, and intended use.
Only then can an oracle and an acceptance rule be selected. The recurring record
shape is
\begin{equation}
  (\text{reference},\text{candidate},\text{metric},
    \text{tolerance},\text{result}).
\end{equation}
The tuple is incomplete unless the reference and candidate carry the metadata
required to make the metric meaningful.

Different claims may inspect the same numerical array. For example, a
Hermiticity residual can verify a software invariant, a comparison with an
independently derived matrix can numerically verify an implementation, and a
comparison with measured band-edge data can contribute to scientific
validation. These are separate interpretations with separate references.

\section{Evidence classes}

\begin{description}
  \item[Software verification.] Establishes that software follows a documented
    contract. Examples include shape and type invariants, exact serialization
    round trips, basis-order checks, command construction, failure behavior,
    and declared symmetry or Hermiticity tests.
  \item[Numerical verification.] Establishes agreement between a numerical
    implementation and independently derived mathematics, or controlled
    convergence toward the stated mathematical problem. Examples include
    analytic limiting cases, independently computed norms, manufactured
    transformations, and stencil-refinement studies.
  \item[Scientific validation.] Evaluates adequacy against independent trusted
    physical, experimental, or scientific reference evidence for a declared
    use. It addresses whether the selected model represents the relevant
    physical behavior, not merely whether its code is internally consistent.
  \item[Uncertainty quantification.] Identifies relevant uncertainty sources and
    characterizes or propagates them to declared outputs. Sensitivity alone is
    not automatically a probabilistic uncertainty statement.
\end{description}

No class implies another. A passing software regression does not establish
numerical convergence. Agreement with the Kohn--Sham parent does not establish
agreement with experiment. A validation comparison does not quantify every
uncertainty source.

\section{Software verification}

Software-verification evidence uses public contracts, schemas, exact language
semantics, or otherwise independently stated behavior as its oracle. For the
present program, relevant subjects include

\begin{itemize}
  \item construction of represented-operator records and their intrinsic
    invariants;
  \item basis, geometry, energy-reference, and spin metadata;
  \item Hermiticity and symmetry analyzers;
  \item compatibility checks that precede arithmetic;
  \item signed difference and residual conventions;
  \item serializer field vocabularies and exact promised round trips;
  \item tight-binding limiting cases and Fourier conventions;
  \item alignment recovery under known synthetic unitary transformations;
  \item continuum-embedding behavior for synthetic envelopes; and
  \item external-tool lifecycle, stream capture, and provenance manifests.
\end{itemize}

Tests should use the supported public surface and an oracle that does not simply
reproduce the production algorithm. Exact contracts use exact acceptance;
numerical contracts use explicit tolerances justified by the operation and
scale. Helpers and fixtures remain distinguishable from evidence-bearing cases.

A round trip establishes only the promised wire representation. A compatibility
result establishes only that the represented prerequisites checked by that
operation agree. Neither proves that recorded provenance is truthful or that the
physical model is adequate.

\section{Numerical verification}

Numerical verification asks whether the implementation approximates the stated
mathematics. Controlled cases are especially important before first-principles
artifacts are available. Examples include

\begin{itemize}
  \item exact recovery of an operator after a known unitary gauge change;
  \item principal angles for analytically constructed subspaces;
  \item residual norms derived independently for small matrices;
  \item Hermitian tight-binding operators in zero-hopping or symmetry-limited
    cases;
  \item Fourier transforms checked against finite hand-constructed lattices;
  \item effective-mass stencils applied to an analytic quadratic dispersion;
  \item continuum solvers checked against solvable potentials where the model
    assumptions match; and
  \item embedding maps checked with synthetic normalized envelopes.
\end{itemize}

Synthetic data are labelled as such. Exact recovery on a controlled example is
strong evidence about the corresponding mathematical implementation, but it is
not a first-principles calculation or physical validation.

\section{Finite-setting convergence}

Suppose a quantity $q_j$ is computed at two successive finite numerical
settings. Observing
\begin{equation}
  |q_{j+1}-q_j|\leq\tau
\end{equation}
establishes stability over those tested settings under the recorded controls. It
does not by itself bound
\begin{equation}
  |q_j-q_{\infty}|.
\end{equation}
An infinite-setting bound requires additional structure, such as a justified
asymptotic error model, monotonicity result, variational bound, or sufficient
higher-resolution evidence. The monograph therefore uses ``stable over the
tested sequence'' unless a stronger result is actually supported.

Cutoff and mesh errors may interact. For an accepted candidate pair
$(E_{\ast},K_{\ast})$ and higher guards $(E_{+},K_{+})$, a finite two-factor
contrast is
\begin{equation}
  I_q=
  \left|
    q(E_{+},K_{+})-q(E_{+},K_{\ast})
    -q(E_{\ast},K_{+})+q(E_{\ast},K_{\ast})
  \right|.
\end{equation}
A small value detects limited interaction in that tested rectangle; it does not
prove global separability of cutoff and sampling errors.

\section{Derived quantities and error amplification}

A pointwise energy tolerance need not control a derivative. For a central
second-difference estimate,
\begin{equation}
  E''(k_0)
  \approx
  \frac{E(k_0+h)-2E(k_0)+E(k_0-h)}{h^2},
\end{equation}
pointwise energy perturbations are amplified by a scale proportional to
$h^{-2}$. Effective-mass verification must therefore vary the stencil or fit
window, report conditioning, identify the selected valley and axes, and check
stability under at least one smaller scale. A converged band plot is not an
effective-mass convergence study.

The same principle applies to subtraction of two large aligned operators,
normalization by a small reference scale, and composition of successive
reduction maps. Numerical cancellation, conditioning, and nonfinite
intermediates must be evaluated at the level of the derived quantity.

\section{Scientific validation}

Scientific validation compares a declared model output with independent
physical evidence appropriate to the intended use. Potential bulk-silicon
references include experimental structural or band-edge measurements and
higher-level theoretical results, but each comparison must match temperature,
geometry, relativistic treatment, observable definition, and uncertainty as far
as the claim requires.

The PBE Kohn--Sham parent is not fitted silently to the experimental silicon gap.
Its disagreement with experiment belongs to parent-model assessment. A
Wannier or tight-binding model is first evaluated against that selected parent
to measure reduction fidelity. If the reduced model is then compared with
experiment, the account must not attribute the parent's discrepancy to the
reduction or claim cancellation as validation without a declared rationale.

For dopants, validation references must match the charge-state interpretation,
periodic or isolated setting, spin and SOC branch, band-edge reference, and
observable definition. A neutral periodic P:Si calculation cannot be validated
as though it were directly the potential of an isolated ionized donor.

\section{Uncertainty quantification}

Relevant uncertainty sources can include

\begin{itemize}
  \item parent-model choice and external-reference uncertainty;
  \item cutoff, mesh, cell-size, and solver discretization;
  \item lattice-fit and valley-location uncertainty;
  \item effective-mass fit-window and stencil sensitivity;
  \item Wannier windows, projections, gauge, and localization choices;
  \item alignment and scalar energy-reference uncertainty;
  \item parameter non-identifiability and optimization variability;
  \item model-class and real-space truncation sensitivity; and
  \item continuum discretization, boundary, and embedding uncertainty.
\end{itemize}

These sources are not combined automatically. A valid propagation statement
requires compatible quantities, an explicit dependence model, and treatment of
correlation where material. Until then, the sources are reported separately as
sensitivity or bounded numerical evidence.

\section{Training, validation, and leakage}

Model fitting uses only the frozen training set. Withheld points, band-edge
quantities, operator blocks, or dopant observables designated for validation
must not influence parameter selection, model-class expansion, stopping rules,
or tolerance adjustment. If a validation failure motivates a new model class,
that class and a new untouched validation design must be recorded explicitly;
the original result remains part of the evidence history.

Visual agreement can support diagnosis but is not an acceptance rule. Plots
must expose the data identity, axes, units, reference, and residual definition,
while machine-readable results retain the exact values used for the decision.

\section{Provenance and reproducibility}

A result is reproducible only relative to retained identities and declared
conditions. The minimum record depends on the operation but generally includes

\begin{itemize}
  \item immutable input and source identities;
  \item code revision and dependency or executable versions;
  \item physical and numerical settings;
  \item exact commands and environment information;
  \item artifact manifests and checksums;
  \item metric and tolerance definitions;
  \item structured outcomes and warnings; and
  \item locations of externally retained large artifacts.
\end{itemize}

Reproduction does not promote the evidence class. Reproducing a tutorial
calculation reproduces tutorial behavior; reproducing a synthetic test
reproduces software or numerical verification; reproducing a production
calculation reproduces that calculation under its recorded conditions.

\section{Current execution-status example}

The repository retains a direct bulk-silicon bootstrap campaign in which nine
SCF and nine linked NSCF Quantum ESPRESSO invocations exited successfully and
reported completion. The owning disposition classifies these observations as
bootstrap scientific-harness development evidence. They are not a canonical
scientific workflow record and do not accept a final cutoff, mesh, lattice
parameter, gap, mass, infinite-basis result, or scientific validation claim.

This example demonstrates why process completion and evidentiary acceptance are
separate. The observations may inform integration and future workflow design
without being discarded or promoted into a production result.

\section{Failure and negative evidence}

A failed acceptance criterion is retained rather than weakened. The diagnosis
may identify an implementation defect, inadequate numerical resolution,
incompatible representation, insufficient model-class expressiveness, or
mismatch between a model and its intended physical use. These dispositions have
different remedies.

Similarly, a well-executed negative result can be scientifically useful: no
model in a frozen hierarchy may satisfy both spectral and operator criteria; no
finite tested crossover radius may satisfy continuum criteria; two reduction
paths may fail to commute within tolerance. The evidence must report the tested
domain and must not generalize beyond it.

\section{Claim-to-evidence discipline}

Later results are included only when each statement can point to an applicable
retained record. Software verification supports contract claims. Numerical
verification supports stated mathematical or convergence claims. Scientific
validation supports adequacy for a declared physical use. UQ supports only the
uncertainties it actually identifies and propagates.

This discipline is not administrative overhead added after computation. It is
how the project prevents a passing test, a successful executable, a smooth plot,
or an aligned matrix from being mistaken for a stronger scientific conclusion.
